\clearpage
\appendix

% ==============================================================================
% Apêndice: Templates de Prompt Aplicados à Questão de Exemplo
% Instância FOLIO id=49 (validação) — Coco Gauff / Roland Garros 2022
% Rótulo original: True
% ==============================================================================

\section{Templates de Prompt Aplicados à Questão de Exemplo}
\label{ap:templates-prompts}

Este apêndice apresenta as nove configurações de prompt utilizadas no experimento,
aplicados a uma mesma instância do conjunto de validação do FOLIO. A instância
selecionada corresponde ao problema sobre tenistas no Roland Garros 2022 (id 49
na partição de validação), cujo rótulo original é \textsc{True}. A utilização
de uma instância comum permite comparar as instruções fornecidas ao modelo em
cada condição sem alterar o conteúdo lógico fundamental do problema. Os prompts
reproduzem literalmente as mensagens enviadas ao modelo durante as inferências,
incluindo o texto de instrução, a ordem dos campos e os rótulos de saída
permitidos. Nos prompts, as premissas em FOL foram normalizadas para ASCII
conforme o pré-processamento adotado no experimento (por exemplo,
$\forall$ foi substituído por \texttt{forall} e $\rightarrow$ por
\texttt{->}). A apresentação introdutória da instância utiliza linguagem natural
para facilitar a leitura, enquanto as listagens preservam a representação ASCII
efetivamente enviada aos modelos.

A instância de referência é composta pelas seguintes premissas e conclusão:

\begin{quote}
\textbf{Premissas (linguagem natural):}
\begin{enumerate}
  \item If a player is ranked highly by the Women's Tennis Association, then
        they are among the most active players in major tennis.
  \item Everyone who lost to Iga Świątek is ranked highly by the Women's
        Tennis Association.
  \item All female tennis players at Roland Garros 2022 lost to Iga Świątek.
  \item Either female tennis players at Roland Garros 2022 or male tennis
        players at Roland Garros 2022.
  \item All male tennis players at Roland Garros 2022 lost to Rafael Nadal.
  \item If Coco Gauff is a player who is ranked highly by the Women's Tennis
        Association or a player who lost to Rafael Nadal, then Coco Gauff is
        not a male tennis player at Roland Garros 2022.
\end{enumerate}
\textbf{Conclusão:} Coco Gauff is among the most active players in major
tennis. \quad \textsc{[True]}
\end{quote}

% ------------------------------------------------------------------------------
\subsection{1. Original}
\label{ap:prompt-original}

Premissas e conclusão apresentadas em lógica de primeira ordem (FOL),
na ordem original do dataset. Rótulo esperado: \textsc{True}.

\begin{lstlisting}[caption={Prompt -- condição Original.},
                   label={lst:prompt-original}]
You are a logic expert.

Given the following premises and a conclusion, determine whether
the conclusion logically follows from the premises.

Premises:
- forall x (RankedHighly(x) -> MostActive(x))
- forall x (LostToSwiatek(x) -> RankedHighly(x))
- forall x (FemalePlayers(x) -> LostToSwiatek(x))
- forall x (FemalePlayers(x) OR MalePlayers(x))
- forall x (MalePlayers(x) -> LostToNadal(x))
- (RankedHighly(cocoGauff) OR LostToNadal(cocoGauff)) -> NOT MalePlayers(cocoGauff)

Conclusion:
MostActive(cocoGauff)

Please reason step by step. Then put your final answer within
\boxed{True}, \boxed{False}, or \boxed{Uncertain}.
\end{lstlisting}

% ------------------------------------------------------------------------------
\subsection{2. Linguagem Natural}
\label{ap:prompt-nl}

Premissas e conclusão apresentadas em inglês, sem notação simbólica.
Rótulo esperado: \textsc{True}.

\begin{lstlisting}[caption={Prompt -- condição Linguagem Natural.},
                   label={lst:prompt-nl}]
You are a logic expert.

Given the following premises and a conclusion, determine whether
the conclusion logically follows from the premises.

Premises:
- If a player is ranked highly by the Women's Tennis Association,
  then they are among the most active players in major tennis.
- Everyone who lost to Iga Swiatek is ranked highly by the Women's
  Tennis Association.
- All female tennis players at Roland Garros 2022 lost to Iga Swiatek.
- Either female tennis players at Roland Garros 2022 or male tennis
  players at Roland Garros 2022.
- All male tennis players at Roland Garros 2022 lost to Rafael Nadal.
- If Coco Gauff is a player who is ranked highly by the Women's Tennis
  Association or a player who lost to Rafael Nadal, then Coco Gauff is
  not a male tennis player at Roland Garros 2022.

Conclusion:
Coco Gauff is among the most active players in major tennis.

Please reason step by step. Then put your final answer within
\boxed{True}, \boxed{False}, or \boxed{Uncertain}.
\end{lstlisting}

% ------------------------------------------------------------------------------
\subsection{3. Premissas Embaralhadas}
\label{ap:prompt-shuffled}

A ordem das premissas foi aleatorizada com semente $42 + \text{id}$.
O conteúdo lógico é idêntico ao da condição Original.
Rótulo esperado: \textsc{True}.

\begin{lstlisting}[caption={Prompt -- condição Premissas Embaralhadas.},
                   label={lst:prompt-shuffled}]
You are a logic expert.

Given the following premises and a conclusion, determine whether
the conclusion logically follows from the premises.

Premises:
- forall x (FemalePlayers(x) OR MalePlayers(x))
- (RankedHighly(cocoGauff) OR LostToNadal(cocoGauff)) -> NOT MalePlayers(cocoGauff)
- forall x (FemalePlayers(x) -> LostToSwiatek(x))
- forall x (LostToSwiatek(x) -> RankedHighly(x))
- forall x (MalePlayers(x) -> LostToNadal(x))
- forall x (RankedHighly(x) -> MostActive(x))

Conclusion:
MostActive(cocoGauff)

Please reason step by step. Then put your final answer within
\boxed{True}, \boxed{False}, or \boxed{Uncertain}.
\end{lstlisting}

% ------------------------------------------------------------------------------
\subsection{4. Premissas Juntas (AND)}
\label{ap:prompt-junto}

Todas as premissas são concatenadas por conjunção lógica em uma única
sentença. Rótulo esperado: \textsc{True}.

\begin{lstlisting}[caption={Prompt -- condição Premissas Juntas.},
                   label={lst:prompt-junto}]
You are a logic expert.

Given the following premises and a conclusion, determine whether
the conclusion logically follows from the premises.

Premises:
(forall x (RankedHighly(x) -> MostActive(x))) AND
(forall x (LostToSwiatek(x) -> RankedHighly(x))) AND
(forall x (FemalePlayers(x) -> LostToSwiatek(x))) AND
(forall x (FemalePlayers(x) OR MalePlayers(x))) AND
(forall x (MalePlayers(x) -> LostToNadal(x))) AND
((RankedHighly(cocoGauff) OR LostToNadal(cocoGauff)) ->
  NOT MalePlayers(cocoGauff))

Conclusion:
MostActive(cocoGauff)

Please reason step by step. Then put your final answer within
\boxed{True}, \boxed{False}, or \boxed{Uncertain}.
\end{lstlisting}

% ------------------------------------------------------------------------------
\subsection{5. Ruído Irrelevante}
\label{ap:prompt-irrelevant}

Uma premissa semanticamente correta mas logicamente irrelevante foi
inserida e as premissas foram reembaralhadas (semente $42 + \text{id} + 1000$).
A premissa inserida neste exemplo foi \texttt{MadeOf(moon, rock)}.
Rótulo esperado: \textsc{True}.

\begin{lstlisting}[caption={Prompt -- condição Ruído Irrelevante.},
                   label={lst:prompt-irrelevant}]
You are a logic expert.

Given the following premises and a conclusion, determine whether
the conclusion logically follows from the premises.

Premises:
- forall x (MalePlayers(x) -> LostToNadal(x))
- forall x (FemalePlayers(x) OR MalePlayers(x))
- forall x (LostToSwiatek(x) -> RankedHighly(x))
- (RankedHighly(cocoGauff) OR LostToNadal(cocoGauff)) -> NOT MalePlayers(cocoGauff)
- forall x (FemalePlayers(x) -> LostToSwiatek(x))
- forall x (RankedHighly(x) -> MostActive(x))
- MadeOf(moon, rock)

Conclusion:
MostActive(cocoGauff)

Please reason step by step. Then put your final answer within
\boxed{True}, \boxed{False}, or \boxed{Uncertain}.
\end{lstlisting}

% ------------------------------------------------------------------------------
\subsection{6. Premissa Faltante}
\label{ap:prompt-missing}

A premissa de índice 4 foi removida --- neste exemplo,
\texttt{forall x (MalePlayers(x) -> LostToNadal(x))} ---, identificada
pelo provador Vampire como necessária para sustentar a conclusão.
Sem ela, a inferência torna-se indeterminada.
Rótulo esperado: \textsc{Uncertain}.

\begin{lstlisting}[caption={Prompt -- condição Premissa Faltante.},
                   label={lst:prompt-missing}]
You are a logic expert.

Given the following premises and a conclusion, determine whether
the conclusion logically follows from the premises.

Premises:
- forall x (RankedHighly(x) -> MostActive(x))
- forall x (LostToSwiatek(x) -> RankedHighly(x))
- forall x (FemalePlayers(x) -> LostToSwiatek(x))
- forall x (FemalePlayers(x) OR MalePlayers(x))
- (RankedHighly(cocoGauff) OR LostToNadal(cocoGauff)) -> NOT MalePlayers(cocoGauff)

Conclusion:
MostActive(cocoGauff)

Please reason step by step. Then put your final answer within
\boxed{True}, \boxed{False}, or \boxed{Uncertain}.
\end{lstlisting}

% ------------------------------------------------------------------------------
\subsection{7. Duplicação de Premissas}
\label{ap:prompt-complex}

Duas premissas existentes foram duplicadas e o conjunto resultante foi
reembaralhado (semente $42 + \text{id} + 2000$), aumentando o tamanho do
contexto sem alterar o conteúdo lógico.
Rótulo esperado: \textsc{True}.

\begin{lstlisting}[caption={Prompt -- condição Duplicação de Premissas.},
                   label={lst:prompt-complex}]
You are a logic expert.

Given the following premises and a conclusion, determine whether
the conclusion logically follows from the premises.

Premises:
- forall x (FemalePlayers(x) -> LostToSwiatek(x))
- forall x (MalePlayers(x) -> LostToNadal(x))
- (RankedHighly(cocoGauff) OR LostToNadal(cocoGauff)) -> NOT MalePlayers(cocoGauff)
- forall x (FemalePlayers(x) OR MalePlayers(x))
- forall x (FemalePlayers(x) -> LostToSwiatek(x))
- forall x (RankedHighly(x) -> MostActive(x))
- forall x (FemalePlayers(x) OR MalePlayers(x))
- forall x (LostToSwiatek(x) -> RankedHighly(x))

Conclusion:
MostActive(cocoGauff)

Please reason step by step. Then put your final answer within
\boxed{True}, \boxed{False}, or \boxed{Uncertain}.
\end{lstlisting}

% ------------------------------------------------------------------------------
\subsection{8. Contradição Injetada}
\label{ap:prompt-contradiction}

A negação de uma premissa existente foi inserida e o conjunto foi
reembaralhado (semente $42 + \text{id} + 3000$). A premissa negada neste
exemplo foi \texttt{forall x (LostToSwiatek(x) -> RankedHighly(x))},
criando uma inconsistência explícita no conjunto de premissas. O critério
de correção nesta condição é o reconhecimento da contradição pelo modelo,
avaliado de forma assistida. Não há um rótulo final esperado nessa condição;
considera-se correto o reconhecimento explícito da incompatibilidade no
raciocínio, independentemente de a resposta final ser \textsc{True},
\textsc{False} ou \textsc{Uncertain}.

\begin{lstlisting}[caption={Prompt -- condição Contradição Injetada.},
                   label={lst:prompt-contradiction}]
You are a logic expert.

Given the following premises and a conclusion, determine whether
the conclusion logically follows from the premises.

Premises:
- forall x (MalePlayers(x) -> LostToNadal(x))
- forall x (LostToSwiatek(x) -> RankedHighly(x))
- forall x (FemalePlayers(x) -> LostToSwiatek(x))
- forall x (FemalePlayers(x) OR MalePlayers(x))
- forall x (RankedHighly(x) -> MostActive(x))
- (RankedHighly(cocoGauff) OR LostToNadal(cocoGauff)) -> NOT MalePlayers(cocoGauff)
- NOT (forall x (LostToSwiatek(x) -> RankedHighly(x)))

Conclusion:
MostActive(cocoGauff)

Please reason step by step. Then put your final answer within
\boxed{True}, \boxed{False}, or \boxed{Uncertain}.
\end{lstlisting}

% ------------------------------------------------------------------------------
\subsection{9. Negação da Conclusão}
\label{ap:prompt-negation}

A conclusão original foi substituída pela sua negação. As premissas
permanecem inalteradas. Como a conclusão original é \textsc{True},
verificar sua negação deve produzir \textsc{False}.
Rótulo esperado: \textsc{False}.

\begin{lstlisting}[caption={Prompt -- condição Negação da Conclusão.},
                   label={lst:prompt-negation}]
You are a logic expert.

Given the following premises and a conclusion, determine whether
the conclusion logically follows from the premises.

Premises:
- forall x (RankedHighly(x) -> MostActive(x))
- forall x (LostToSwiatek(x) -> RankedHighly(x))
- forall x (FemalePlayers(x) -> LostToSwiatek(x))
- forall x (FemalePlayers(x) OR MalePlayers(x))
- forall x (MalePlayers(x) -> LostToNadal(x))
- (RankedHighly(cocoGauff) OR LostToNadal(cocoGauff)) -> NOT MalePlayers(cocoGauff)

Conclusion:
NOT (MostActive(cocoGauff))

Please reason step by step. Then put your final answer within
\boxed{True}, \boxed{False}, or \boxed{Uncertain}.
\end{lstlisting}
